\chapter{Les pièges}

Monte-Carlo a une particularité désagréable : quand il est faux, il ne plante
pas. Il rend un nombre plausible. Ce chapitre rassemble les erreurs qui coûtent
une journée, et les signes qui les trahissent.

\section{Le générateur aléatoire}

Toute la méthode repose sur la qualité du hasard. Un générateur médiocre
introduit des \textbf{corrélations} entre échantillons — et des échantillons
corrélés valent moins que leur nombre.

\begin{nkpiege}[N'utilisez jamais \texttt{rand()} de la bibliothèque C]
Trois défauts, tous rédhibitoires ici :
\begin{itemize}
  \item \textbf{période courte} : sur certaines implémentations $2^{31}$, ce
        qui est atteint en quelques secondes de rendu ;
  \item \textbf{bits de poids faible quasi périodiques} : \texttt{rand() \% 2}
        peut alterner régulièrement. Tout code qui prend un modulo est
        contaminé ;
  \item \textbf{état global} : inutilisable en multi-fils sans verrou, et un
        verrou sur le générateur détruit la parallélisation.
\end{itemize}
Le dépôt fournit mieux : \texttt{math::NkRand} et le XORShift64* de
\nkfile{NKContainers/String/NkStringUtils.cpp}.
\end{nkpiege}

\begin{nkpiege}[Une graine par fil d'exécution, sans exception]
Si deux fils partagent la même graine, ils tirent la \textbf{même séquence} :
vous ne multipliez pas vos échantillons, vous dupliquez le même. Le bruit ne
diminue plus, et rien ne l'annonce.

La bonne pratique tient en une ligne : semer chaque fil avec un identifiant
distinct — le numéro de ligne d'image, l'indice du fil, l'indice du pixel —
mélangé par une fonction de hachage. Un générateur par fil, jamais partagé.
\end{nkpiege}

\begin{nkcode}{Un generateur par fil, seme sans collision}
// Melangeur de bits (finaliseur de MurmurHash3) : deux graines voisines
// donnent des sequences totalement decorrelees. Semer avec `indice` brut
// produirait au contraire des sequences qui demarrent presque pareil.
static inline nk_uint64 NkMelanger(nk_uint64 x) {
    x ^= x >> 33;  x *= 0xFF51AFD7ED558CCDull;
    x ^= x >> 33;  x *= 0xC4CEB9FE1A85EC53ull;
    x ^= x >> 33;
    return x;
}

class NkRngLocal {
    public:
        explicit NkRngLocal(nk_uint64 indice, nk_uint64 graineGlobale = 0x9E3779B9ull)
            : mEtat(NkMelanger(indice + graineGlobale)) {
            if (mEtat == 0)
                mEtat = 0x853C49E6748FEA9Bull;   // XORShift meurt sur zero
        }

        nk_uint64 Suivant() {                    // XORShift64*
            mEtat ^= mEtat >> 12;
            mEtat ^= mEtat << 25;
            mEtat ^= mEtat >> 27;
            return mEtat * 0x2545F4914F6CDD1Dull;
        }

        // [0,1) en float32 : on prend les 24 bits de POIDS FORT. Les bits de
        // poids faible sont les moins bons de tous les generateurs rapides.
        float32 Flottant() {
            return (float32)(Suivant() >> 40) * (1.f / 16777216.f);
        }

    private:
        nk_uint64 mEtat;
};
\end{nkcode}

\section{Le biais silencieux}

Rappel du chapitre 2 : un estimateur biaisé est \emph{stable et faux}. Il ne
tremble pas, il ment — et l'augmentation du nombre d'échantillons ne fait que
resserrer l'estimation \emph{autour de la mauvaise valeur}.

\begin{center}
\begin{tabular}{@{}p{5.2cm}p{4.2cm}p{4.4cm}@{}}
\toprule
\textbf{Cause} & \textbf{Symptôme} & \textbf{Correction} \\
\midrule
oubli de la division par $p$ & résultat faux d'un facteur constant & diviser, toujours \\
coupure à profondeur fixe & image trop sombre & roulette russe \\
densité mal normalisée & décalage proportionnel & vérifier que $\int p = 1$ \\
région jamais échantillonnée & contribution absente & garantir $p > 0$ où $f \neq 0$ \\
\bottomrule
\end{tabular}
\end{center}

\begin{nkretenir}[Le test qui démasque un biais]
Faites tourner votre estimateur avec $N$, puis $100N$ échantillons.

\begin{itemize}
  \item si l'écart entre les deux résultats \textbf{rétrécit d'un facteur 10}
        (racine de 100), c'est du bruit : tout va bien ;
  \item si les deux résultats sont \textbf{proches l'un de l'autre mais
        différents de la valeur attendue}, c'est un biais. Aucun temps de
        calcul ne le corrigera.
\end{itemize}
Testez toujours sur un cas dont vous connaissez la réponse : une sphère
uniformément éclairée, une intégrale calculable à la main, un jeu dont l'issue
optimale est connue.
\end{nkretenir}

\section{L'accumulation numérique}

Additionner un million de \texttt{float32} en série perd les petites
contributions : quand la somme devient grande, l'ajout d'une petite valeur est
absorbé par l'arrondi.

\begin{nkpiege}[Accumuler en \texttt{float64}, ou par sommation compensée]
Le remède le moins cher est d'accumuler en \texttt{float64} et de ne convertir
qu'à la fin. Si la mémoire l'interdit (un buffer d'accumulation d'image en
haute définition), employez la sommation de Kahan, qui garde trace de l'erreur
d'arrondi et la réinjecte.

Symptôme typique quand on l'ignore : une image qui \textbf{cesse de converger}
au-delà de quelques milliers d'échantillons par pixel — le bruit ne descend
plus, alors que la théorie dit qu'il devrait.
\end{nkpiege}

\section{Les \emph{fireflies}}

Des pixels blancs isolés sur une image sombre. Cause : un échantillon a touché
une contribution énorme avec une densité minuscule, donc $f/p$ explose.

\begin{itemize}
  \item \textbf{mauvaise solution} : plafonner brutalement les valeurs
        (\emph{clamping}). Ça marche visuellement, mais ça \textbf{introduit un
        biais} : on retire de l'énergie qui existait vraiment ;
  \item \textbf{bonne solution} : mieux échantillonner ce qui les cause —
        tirage direct des lumières, échantillonnage à importances multiples.
        On s'attaque à la densité, pas au symptôme.
\end{itemize}

\section{La parallélisation}

Monte-Carlo se parallélise magnifiquement : les échantillons sont indépendants.
Deux réserves seulement.

\begin{itemize}
  \item \textbf{le déterminisme}. Une image rendue sur 8 fils ne sera pas
        identique au bit près à la même rendue sur 4 fils, si l'ordre
        d'accumulation change. Pour un rendu reproductible — indispensable
        pour comparer deux versions du moteur — il faut que chaque
        \textbf{échantillon} ait sa graine propre, dérivée de son indice et non
        de l'ordre d'exécution ;
  \item \textbf{le faux partage} (\emph{false sharing}). Si deux fils
        accumulent dans des cases voisines du même tableau, ils se disputent la
        même ligne de cache et le programme ralentit sans raison apparente.
        Donnez à chaque fil son accumulateur, aligné, et fusionnez à la fin.
\end{itemize}

\section{Récapitulatif : la liste à cocher}

\begin{nkretenir}[Avant de croire un résultat Monte-Carlo]
\begin{enumerate}
  \item Est-ce que je divise par la densité de tirage, partout ?
  \item Ma densité est-elle strictement positive partout où la fonction ne
        l'est pas ?
  \item Ai-je testé sur un cas dont je connais la réponse exacte ?
  \item L'erreur décroît-elle bien en $1/\sqrt{N}$ quand je multiplie les
        échantillons ?
  \item Chaque fil a-t-il sa propre graine ?
  \item Est-ce que j'accumule dans un type assez large ?
  \item Est-ce que je coupe quelque chose « pour aller plus vite » sans
        compenser ?
\end{enumerate}
Sept questions. Elles coûtent dix minutes et font gagner des journées.
\end{nkretenir}

\section{Pour aller plus loin}

\begin{itemize}
  \item \textbf{Suites à faible discordance} (Halton, Sobol) : la
        stratification qui survit à la grande dimension ;
  \item \textbf{échantillonnage à importances multiples} : combiner
        proprement plusieurs densités, chacune bonne pour une partie du
        problème ;
  \item \textbf{Metropolis-Hastings et MCMC} : quand la bonne densité est
        impossible à échantillonner directement, on construit une marche
        aléatoire qui l'explore. C'est la suite naturelle de ce cours ;
  \item \textbf{estimateurs à variance réduite} : variables de contrôle,
        variables antithétiques.
\end{itemize}
